\scanpage{50}

$\unsure{c'|c})\Longrightarrow b'$\quad Soit en effet $\Gamma$ un groupe comme dans b'), supposons d'abord $l\ne p$.
Soit $A=\bar k(\Gamma)$ l'algèbre \inserted{de $\Gamma$} à coefficients dans la clôture algébrique $\bar k$ de $k$, \inserted{$\unsure{\mathfrak g}_k$ opérant de façon évidente.} c'est une algèbre déf/k, commutative et \deleted{semi-simple,} \uline{séparable} (car $[\Gamma:e]$ est premier à $p$), donc en vertu de c') on a $H^i(k,G)=0$ si $i>n$, $G$ désignant le groupe des éléments inversibles de $A$.
\figure{\begin{tikzpicture}[x=1cm,y=1cm]\draw (0,0)--(-.05,.4)..controls (1,.62) and (3.7,.4)..(4.7,.6)..controls (4.9,.65) and (4.9,.7)..(4.9,.78);\end{tikzpicture}}
Soit $G_1$ le sous-groupe de $G$ formé des $x\in G$ tels que $\unsure{N}_{A/\bar k}x=1$ (où \deleted{dét $\unsure{=}$} $N_{A/k}x={}$dét. de la translation op. par $x$ dans $A$), c'est un sous-groupe \uline{défini sur $k$} (noyau de l'homomorphisme $N_{A/k}:G\longrightarrow\bar k^*$, qui est déf/k).
Comme $\Gamma$ est \uline{simple} comme $\unsure{\mathfrak g}$-module, \inserted{et que $\Gamma\subset G$} il en résulte que \deleted{$\unsure{N_{A/k}}$} l'on a $\Gamma\subset G_1$, ou $\Gamma\cap G_1=\deleted{\unsure{\{0\}}}\{e\}$, et on vérifie facilement que ce deuxième cas est exclu (distinguer le cas \deleted{\unsure{pair}} \inserted{$l$} impair, et le cas $l=2$ \deleted{$\unsure{;} $}, on note dans ce dernier cas qu'on dans ce cas a $\unsure{N}_{A/k}\xi=\pm1$ si $\xi\in\Gamma$).
\figure{\begin{tikzpicture}[x=1cm,y=1cm]\draw (0,0)..controls (0,.15) and (-.04,.35)..(-.04,.4)..controls (1,.6) and (2.9,.4)..(3.3,.6);\end{tikzpicture}}
Donc $\Gamma\subset G_1$. Comme $l\Gamma\deleted{\unsure{=0}}=\{0\}$, il en résulte que $\Gamma$ est contenu dans le noyau de l'hom. $G_1\xrightarrow{l}G_1$.
Ce noyau est d'ailleurs d'ordre $l^{n_1}$, où $n_1=n-1$ est la dim de $G_1$, \unsure{donc} $n$ l'ordre de $\unsure{G|\Gamma}$.

\mspage{49}

